某大型医药流通企业在全国范围内运营着覆盖广泛的物流网络，包括5个枢纽仓库、41家省级仓库和200余家地市级仓库，自有车辆超过2000辆（其中冷藏车400余辆），终端配送覆盖100000余家药房和10000余家医院。在药品配送过程中，冷链药品对运输工具和时效有严格要求——除普通药品可使用常温车外，其余药品均需使用冷藏车运输，且送货时限通常允许3至4天，最长不超过一周。

本次竞赛围绕该企业的冷链药品运输优化展开，包含三个递进问题：

\textbf{问题一：运输费用分析。} 根据附件1提供的车辆运营成本与承载能力参数，以及附件2给出的南通市崇川区某医药企业仓库出发的冷链药品排货信息，建立运输成本核算模型，分析计算相应的运输总费用。

\textbf{问题二：拼车优化运输。} 在问题一的基础上，通过拼车（即多批货物合并运输）降低运营成本。装卸货工作时间限制为每日9:00至17:00，每次装卸约需2小时。假设所有订单信息可提前一周预知，设计改进的拼单运输和车辆调度方案，实现运输成本的最小化。

\textbf{问题三：派车单评估与优化。} 根据附件3提供的某一周终端运单需求数据和附件4给出的仓库至终端派车单数据，建立合理的评估指标体系，评价当前派车单安排的合理性，并在保证取送货时限的前提下提出优化方案。